\chapter[Materiales y métodos]{Materiales y métodos}
\label{cp:metodos}
% Extensión estimada: 20--25 páginas.
% Fuentes: sección "Configuración experimental" y subsecciones "Diseño del experimento" del documento de
% experimentación, Propuesta (datasets, técnicas), planes en Codigo/docs/ideas/ (en particular
% split_lesion_semillas_plan_20260912.md). Incorpora el protocolo unificado pedido en el feedback del tutor
% del 12/09/2026 (Codigo/docs/feedback_20260912.md).
%
% Feedback del tutor (14/09/2026):
%  - "Particiones y auditoría de fuga de información" se renombra: el término no se entiende fuera de
%    contexto. Ahora la subsección se llama por lo que produce (la partición disjunta por lesión) y el
%    término "fuga de información (data leakage)" se introduce dentro del texto, definido.
%  - La subsección "Métricas" repetía el marco teórico. Se fusiona con la de selección: el marco teórico
%    define qué es cada métrica, aquí se justifica cuál se usa para decidir.
%  - "Robustez, calibración e interpretabilidad" no decía a qué se refería; se renombra por extenso.
%  - "Sistema propuesto" pasa a llamarse "pipeline", que es el término corriente en el área para la
%    arquitectura completa: clasificador, kNN y generador encadenados.
%  - Se incorpora aquí "Uso de un VLM generativo como clasificador", que estaba en el marco teórico pese a
%    ser parte del enfoque propio.
%
% Revisión 25/09/2026: recibe del marco teórico la prosa ya escrita sobre las categorías de HAM10000
% (§HAM10000), los costos asimétricos de confusión (§selección de métricas) y las partes del reporte
% (§generación de reportes). Ver fichas/02-marco-teorico.md, decisión 14.

\pendiente{Este capítulo describe todo lo necesario para reproducir los experimentos: datos, particiones,
modelos, \emph{pipeline} y protocolo de evaluación. Los resultados se presentan en los capítulos
siguientes.}

\section{Metodología de trabajo}
\label{sec:met_metodologia}
\pendiente{Se describe el ciclo iterativo \emph{hipótesis $\rightarrow$ ejecución $\rightarrow$ análisis
$\rightarrow$ decisión} con registro sistemático: un plan con hipótesis y criterio de éxito, la ejecución y
un resultado con veredicto binario. Varias etapas del trabajo nacieron de refutar la hipótesis anterior, de
modo que el orden de los capítulos de resultados sigue el orden en que se tomaron las decisiones: baseline
supervisado, VLMs en clasificación, refutación de la hipótesis del desbalance, especialización de dominio,
generación de reportes y ensamble en cascada, y análisis de rigor y robustez.}

\section{Conjuntos de datos y particiones}
\label{sec:met_datasets}

\subsection{HAM10000}
\label{sec:met_ham}

El clasificador asigna cada imagen a una etiqueta de un inventario cerrado, y ese inventario es el de
HAM10000, con siete categorías: \gls{mel}, \gls{nv}, \gls{bcc}, \gls{akiec}, \gls{bkl}, \gls{df} y
\gls{vasc}. Según su gravedad, NV, BKL, DF y VASC son benignas, AKIEC reúne lesiones premalignas y
carcinomas in situ, y MEL y BCC son malignas~\cite{tschandl2018ham10000}. Esa gradación determina
cuánto cuesta cada confusión del clasificador, y con ello qué métricas se usan para juzgarlo
(Sección~\ref{sec:met_seleccion}).

\pendiente{Se presentan la composición del conjunto, la distribución de clases y su desbalance.}

\subsection{Partición disjunta por lesión}
\label{sec:met_particiones}
\pendiente{HAM10000 contiene varias imágenes de una misma lesión y la partición original se estratificó por
imagen, de modo que imágenes de una misma lesión quedaron repartidas entre entrenamiento y test. Esta
subsección explica ese problema, conocido como fuga de información entre particiones (\emph{data leakage}),
lo cuantifica (el 36\% del test comparte lesión con entrenamiento; el 63\% en melanoma) y describe cómo se
midió su efecto sobre las métricas: reevaluación sobre el subconjunto de test cuya lesión no aparece en
entrenamiento ni en validación, y control con un modelo que nunca vio HAM10000 (DermLIP \emph{zero-shot}).

Se describe luego la partición disjunta por lesión que se construyó como reemplazo, con test idéntico al
original y validación agrupada (\texttt{StratifiedGroupKFold}), sus tamaños finales y qué resultados se
recalcularon sobre ella.}

\subsection{SkinCaRe, MM-Skin y HIBA}
\pendiente{Se describe el rol de cada conjunto: entrenamiento de la generación (SkinCAP y SkinCoT),
evaluación con referencia (MM-Skin) y evaluación fuera de distribución (HIBA, solo test). Se indica la
modalidad de cada uno, en particular que SkinCaRe es de fotografía clínica y HAM10000 dermatoscópico, con
remisión a la Sección~\ref{sec:marco_generalizacion}.}

\subsection{Preprocesamiento y aumento de datos}
\pendiente{Se detallan la resolución, la normalización, los aumentos de entrenamiento y el TTA.}

\section{Modelos evaluados}
\label{sec:met_modelos}

\subsection{Clasificadores supervisados}
\pendiente{Se describen DenseNet121, ViT-B/16, Swin-Tiny y ConvNeXt-Tiny: pre-entrenamiento, cabeza de
clasificación, pérdidas y estrategias de desbalance evaluadas.}

\subsection{Modelos visión--lenguaje médicos generales}
\pendiente{Se presentan Mini-VLM (Phi-3.5-Vision) y LLaVA-Med: arquitectura, tamaño y motivo de su
elección.}

\subsection{Modelo visión--lenguaje dermatológico: DermLIP}
\pendiente{Se describen la arquitectura y el pre-entrenamiento de DermLIP (Derm1M) y sus modos de uso:
\emph{zero-shot}, \emph{linear probe}, kNN y SVM sobre \emph{embeddings}.}

\subsection{Generador dermatológico: SkinVL}
\pendiente{Se presenta SkinVL y el motivo de incluirlo como generador especializado de comparación.}

\section{Clasificación con modelos visión--lenguaje}
\label{sec:met_clasif_vlm}
\pendiente{Obtener una clase de un modelo generativo admite varias formas, y la elección condiciona el
resultado. Esta sección presenta las cuatro estrategias evaluadas —respuesta libre (\emph{zero-shot}),
\emph{likelihood scoring} sobre las etiquetas candidatas, razonamiento estructurado (\emph{Chain-of-Thought})
y ajuste fino con LoRA— junto con los prompts, el parseo de respuestas y la configuración de cada una.}

\section{Generación de reportes}
\label{sec:met_generacion}

El reporte dermatológico deja por escrito la lectura de la lesión, y de sus partes el trabajo se ocupa
de dos: la descripción de lo que se ve y la impresión diagnóstica. Son las que corresponden a las dos
tareas del problema, y las que los conjuntos de pares imagen--texto del área anotan junto a cada
imagen~\cite{shen2025skincare}. El diagnóstico diferencial y la recomendación de manejo, que un reporte
clínico completo también lleva, quedan fuera del alcance.

\pendiente{Se describen el ajuste fino del generador sobre SkinCaRe con LoRA, el formato concreto del
reporte generado y la configuración de decodificación.}

\section{Pipeline propuesto: ensamble en cascada consciente de la incertidumbre}
\label{sec:met_pipeline}

\subsection{Cascada clasificador--generador (\emph{grounding})}
\pendiente{Se explica por qué separar el diagnóstico de la redacción, cómo se inyecta el diagnóstico en el
prompt y qué variantes se evalúan (clasificador, DermLIP y \emph{oracle} como cota superior). Un diagrama
resume el pipeline.}

\subsection{Grounding por recuperación}
\pendiente{Se describen la búsqueda de vecinos visuales con DermLIP (kNN, $k=3$) y la construcción del
contexto.}

\subsection{Ruteo por incertidumbre en tres niveles}
\pendiente{Se explica cómo el acuerdo entre clasificador y kNN determina el nivel de cada caso (afirmativo,
diferencial o abstención con derivación) y se presentan las plantillas de cada nivel.}

\section{Protocolo de evaluación}
\label{sec:met_protocolo}

\subsection{Selección de modelos, métricas y criterio de decisión}
\label{sec:met_seleccion}

Las confusiones entre categorías no tienen el mismo costo clínico. Confundir MEL con NV retrasa la
extirpación de una lesión cuya letalidad depende del estadio en que se diagnostica; confundir NV con MEL
lleva a extirpar una lesión benigna sin necesidad; las confusiones entre las clases benignas tienen menor
consecuencia~\cite{tschandl2018ham10000,nci2026melanoma}. Esa asimetría obliga a medir el desempeño clase
por clase, además de en conjunto, y a dar a las clases minoritarias el mismo peso que a la mayoritaria.

\pendiente{El Capítulo~\ref{cp:marco} define qué mide cada métrica; esta subsección indica cuáles se usan
para decidir y por qué. Se justifican la sensibilidad macro como métrica principal de clasificación y el F1
macro como complemento, el umbral de especificidad $\geq 0.85$ fijado en validación para la detección de
melanoma, las métricas de calibración y las dos de generación (Nivel 1 con referencia; Nivel 2 con
\emph{clinical entity recall} y \emph{diagnostic consistency}), con la distinción entre fidelidad diagnóstica
y validez clínica.

Se describen a continuación los barridos de hiperparámetros y el criterio de selección, que opera
enteramente en validación: el test no interviene en ninguna decisión. Hay que distinguir dos niveles. El
\emph{checkpoint} de cada corrida es la época con mejor F1 ponderado en validación. La corrida ganadora
entre las del barrido, en cambio, se elige por \emph{F1 macro} en validación, porque el F1 ponderado está
dominado por la clase mayoritaria ---NV es el 67\,\% de las imágenes--- y la sensibilidad macro, que sería
la alternativa natural por ser la métrica principal de este trabajo, ignora los falsos positivos y por lo
tanto premia a los modelos que sobre-predicen las clases minoritarias. El F1 macro es el único de los tres
que pondera las clases por igual y a la vez penaliza los falsos positivos. La sensibilidad macro se reporta
como resultado principal, pero no se usa para elegir.

Los resultados se reportan como media $\pm$ desvío entre semillas.}

\subsection{Protocolos de inferencia}
\pendiente{Una tabla única indica, para cada resultado reportado, si se usa vista única o TTA, qué
\emph{checkpoint} y qué partición.}

\subsection{Análisis estadístico}
\pendiente{Se indica qué intervalo o prueba corresponde a cada métrica y comparación: \emph{bootstrap}
pareado para diferencias de métricas agregadas, McNemar para aciertos por imagen y Wilcoxon para
\emph{recall} por clase.}

\subsection{Análisis complementarios: validación externa, perturbaciones, calibración e interpretabilidad}
\pendiente{Además de las métricas sobre el test, se realizan cuatro análisis que responden preguntas
distintas de ``cuánto acierta el modelo'': si el desempeño se sostiene en otra población (evaluación en
HIBA), si se sostiene bajo peores condiciones de captura (perturbaciones de adquisición: ruido, desenfoque,
compresión JPEG y pérdida de resolución), si las probabilidades que entrega el modelo son confiables
(calibración post-hoc ajustada en validación) y si el modelo mira la lesión o el contexto de la imagen
(Grad-CAM). Esta subsección describe el procedimiento de cada uno.}

\section{Infraestructura y reproducibilidad}
\label{sec:met_infra}
\pendiente{Se describen los recursos de cómputo (GPU local RTX 4070, clúster TOKO con RTX 4080 y Mendieta con
A30), el uso de SLURM, el seguimiento de experimentos y la organización del código y del registro de
ejecuciones. El detalle está en el Apéndice~\ref{ap:reproducibilidad}.}
